\documentclass[10pt,a4paper]{amsart}
\usepackage[margin=19mm]{geometry}
\usepackage{amsmath,amssymb,mathtools}
\usepackage[scr=rsfs,cal=euler]{mathalfa}
\usepackage{xfrac}
\usepackage[unicode,hidelinks,bookmarks=false]{hyperref}
\newcommand{\ie}{i.e.\ }
\newcommand{\cf}{cf.\ }
\newcommand{\resp}{respectively\ }
\newcommand{\Subsection}{\subsection}
\providecommand{\nobreakdash}{\nobreak\hbox{-}}
\setlength{\emergencystretch}{2em}
\multlinegap0pt
\usepackage{bm}
\usepackage{bbm}
\usepackage{dsfont}
\usepackage{xspace}
\usepackage{cancel}
\usepackage{mathtools}
\usepackage{amsthm}
\makeatletter
\@ifundefined{theorem}{\newtheorem{theorem}{Theorem}}{}
\@ifundefined{assumption}{\newtheorem{assumption}[theorem]{Assumption}}{}
\@ifundefined{lemma}{\newtheorem{lemma}[theorem]{Lemma}}{}
\makeatother
\newcommand{\E}{\mathbb{E}}
\newcommand{\ip}[2]{\left\langle #1,#2\right\rangle}
\newcommand{\norm}[1]{\left\lVert #1\right\rVert}
\title[Nonlinear Two-Time-Scale Stochastic Approximation: A Sharp P]{Nonlinear Two-Time-Scale Stochastic Approximation: A Sharp Phase Transition and How to Beat It}
\author{Dhruv Sarkar, Vaneet Aggarwal}
\date{Source: arXiv:2606.14488v1 (CC BY 4.0)}
\begin{document}
\maketitle

\noindent\emph{Source and licence.} Nonlinear Two-Time-Scale Stochastic Approximation: A Sharp Phase Transition and How to Beat It. Version arXiv:2606.14488v1 (CC BY 4.0), distributed under the Creative Commons Attribution 4.0 International licence, as recorded in the arXiv licence metadata for this exact version and shown on the abstract page https://arxiv.org/abs/2606.14488v1. Full rights notice: licenses/arxiv-2606.14488-CC-BY-4.0.txt.\par\smallskip

\noindent\textit{Editorial scope.} Counted unit: the Lemma whose source label is \texttt{lem:fast-moments} in the article (source0.tex lines 1039--1049, source label \texttt{lem:fast-moments}), delivered together with the article's own complete original proof of it (source0.tex lines 1051--1197, one \texttt{proof} environment of 147 lines). No mathematics is written, repaired or re-derived: statement and proof are the article's own text line for line. Required context retained uncounted: ass:stability (lines 232--252); eq:moment-noise (lines 246--250). Every definition, display and label of those lines is retained unchanged. Omitted from this package and not redistributed: 1--231, 251--1038, 1050--1050, 1198--2132. Nothing inside a retained excerpt is omitted or altered: every retained body is delivered exactly as the article's lines are written, so no omission receipt of any kind accompanies this package. Citations: the retained excerpts carry no citation command at all, so no bibliography entry of the article is transcribed and no reference is redistributed. Reference adaptation: none was needed and none was made; every reference inside the delivered unit resolves inside it, so no unresolved reference or citation is produced. Printed slips kept literal: no printed word, symbol, hypothesis or number of the counted statement or of its proof was altered. Layout and class: the article's own class is \texttt{article} with packages including arxiv, inputenc, fontenc, amsmath, amssymb, amsthm, mathtools, bbm, booktabs, enumitem, microtype, hyperref; the wrapper is the lane's pinned \texttt{amsart} editorial wrapper at 10pt on A4, which restates the theorem-like environments the retained text needs and adds the article's own macro definitions that the retained text needs (\texttt{\textbackslash E}, \texttt{\textbackslash ip}, \texttt{\textbackslash norm}), with extra wrapper packages (bm, bbm, dsfont, xspace, cancel, mathtools). The delivered numbering is the unit's own. Assets: no retained span contains a figure environment, a graphics-inclusion command, a PDF-page inclusion, a table or a TikZ picture; no image file is copied into this package, the package declares no asset dependency and no image file is redistributed with it. Licence: version arXiv:2606.14488v1 of this article is distributed under the Creative Commons Attribution 4.0 International licence, as recorded in the arXiv licence metadata for this exact version and shown on the abstract page https://arxiv.org/abs/2606.14488v1. A full rights notice is delivered at licenses/arxiv-2606.14488-CC-BY-4.0.txt. Characters: every retained excerpt is pure ASCII, so the delivered engine, XeLaTeX with its default Latin Modern text font, reports no missing character.
\begin{assumption}[Stability and noise]\label{ass:stability}
There exist constants $\lambda,b>0$ such that, for all $k\ge0$,
\begin{equation}
  \norm{I-\alpha_k A}\le 1-\lambda\alpha_k,
  \qquad
  \norm{I-\beta_k B}\le 1-b\beta_k. \label{eq:matrix-contraction}
\end{equation}
The matrix $A$ is invertible. The noise sequences are martingale differences:
\begin{equation}
  \E[\xi_{k+1}\mid\mathcal F_k]=0,
  \qquad
  \E[\zeta_{k+1}\mid\mathcal F_k]=0. \label{eq:mds}
\end{equation}
Moreover, for a finite constant $\sigma_4$,
\begin{equation}
  \E[\norm{\xi_{k+1}}^4\mid\mathcal F_k]\le \sigma_4^4,
  \qquad
  \E[\norm{\zeta_{k+1}}^2\mid\mathcal F_k]\le \sigma_4^2. \label{eq:moment-noise}
\end{equation}
The initial conditions $X_0,Y_0$ are deterministic.
\end{assumption}

\begin{lemma}[Fast moments]\label{lem:fast-moments}
Under Assumption \ref{ass:stability}, there exists $C<\infty$ such that, for every $k\ge0$,
\begin{equation}
  \E\norm{X_k}^2\le C\alpha_k,\qquad
  \E\norm{X_k}^4\le C\alpha_k^2. \label{eq:fast-second-fourth}
\end{equation}
Consequently, for every $\rho\in[0,1]$,
\begin{equation}
  \E\norm{X_k}^{2+2\rho}\le C\alpha_k^{1+\rho}. \label{eq:fast-interp}
\end{equation}
\end{lemma}

\begin{proof}
Let $D_k=I-\alpha_kA$. By Assumption \ref{ass:stability}, $\norm{D_k}\le1-\lambda\alpha_k$.

\paragraph{Second moment.}
Conditioning on $\mathcal F_k$ and using $\E[\xi_{k+1}\mid\mathcal F_k]=0$,
\[
\begin{aligned}
  \E[\norm{X_{k+1}}^2\mid\mathcal F_k]
  &=\E[\norm{D_kX_k+\alpha_k\xi_{k+1}}^2\mid\mathcal F_k]\\
  &=\norm{D_kX_k}^2
    +2\alpha_k\ip{D_kX_k}{\E[\xi_{k+1}\mid\mathcal F_k]}
    +\alpha_k^2\E[\norm{\xi_{k+1}}^2\mid\mathcal F_k]\\
  &\le (1-\lambda\alpha_k)^2\norm{X_k}^2+\sigma_4^2\alpha_k^2.
\end{aligned}
\]
Since $\lambda\alpha_k\le1$ after increasing $k_0$ if necessary, $(1-\lambda\alpha_k)^2\le1-\lambda\alpha_k$. Thus
\begin{equation}
  m_{k+1}\le (1-\lambda\alpha_k)m_k+\sigma_4^2\alpha_k^2,\qquad
  m_k\coloneqq \E\norm{X_k}^2. \label{eq:m-recursion}
\end{equation}
We prove $m_k\le C_2\alpha_k$ by induction. Since $X_0$ is deterministic and $\alpha_0$ in the display below denotes the fixed numerator in the stepsize formula, choose $C_2$ so large that $m_0\le C_2\alpha_0 k_0^{-a}$. Because $\alpha_k=\alpha_0\bar k^{-a}$,
\[
  \frac{\alpha_{k+1}}{\alpha_k}=\left(\frac{\bar k}{\bar k+1}\right)^a
  =\left(1+\frac1{\bar k}\right)^{-a}.
\]
The inequality $(1+t)^{-a}\ge1-at$ for $t\ge0$ gives
\begin{equation}
  \alpha_{k+1}\ge \alpha_k\left(1-\frac{a}{\bar k}\right). \label{eq:alpha-lower}
\end{equation}
Since $a<1$, we have $\bar k^{-1}=o(\alpha_k)$ because $\alpha_k=\alpha_0\bar k^{-a}$. Hence, after increasing $k_0$ if necessary,
\begin{equation}
  \frac{a}{\bar k}\le \frac{\lambda}{2}\alpha_k,
  \qquad k\ge0. \label{eq:k0-condition-second}
\end{equation}
Combining \eqref{eq:alpha-lower} and \eqref{eq:k0-condition-second},
\begin{equation}
  \alpha_{k+1}\ge \alpha_k\left(1-\frac{\lambda}{2}\alpha_k\right). \label{eq:alpha-rec-compare}
\end{equation}
Assume $m_k\le C_2\alpha_k$. From \eqref{eq:m-recursion},
\[
  m_{k+1}\le C_2\alpha_k(1-\lambda\alpha_k)+\sigma_4^2\alpha_k^2
  =C_2\alpha_k\left(1-\lambda\alpha_k+\frac{\sigma_4^2}{C_2}\alpha_k\right).
\]
Choose $C_2\ge 2\sigma_4^2/\lambda$. Then
\[
  m_{k+1}\le C_2\alpha_k\left(1-\frac{\lambda}{2}\alpha_k\right)
  \le C_2\alpha_{k+1},
\]
where the last inequality uses \eqref{eq:alpha-rec-compare}. This completes the induction and proves $m_k\le C_2\alpha_k$.

\paragraph{Fourth moment.}
Let $u=D_kX_k$ and $v=\alpha_k\xi_{k+1}$. The identity
\[
  \norm{u+v}^4=(\norm{u}^2+2\ip{u}{v}+\norm{v}^2)^2
\]
expands as
\begin{align*}
  \norm{u+v}^4
  &=\norm{u}^4+4\norm{u}^2\ip{u}{v}
    +2\norm{u}^2\norm{v}^2+4\ip{u}{v}^2\\
  &\qquad +4\ip{u}{v}\norm{v}^2+\norm{v}^4.
\end{align*}
Taking conditional expectation, the term $4\norm{u}^2\ip{u}{v}$ vanishes because $\E[v\mid\mathcal F_k]=0$. For the remaining terms, using $\norm{u}\le\norm{X_k}$ and \eqref{eq:moment-noise},
\begin{align*}
  \E[2\norm{u}^2\norm{v}^2\mid\mathcal F_k]
  &\le 2\sigma_4^2\alpha_k^2\norm{X_k}^2,\\
  \E[4\ip{u}{v}^2\mid\mathcal F_k]
  &\le 4\norm{u}^2\E[\norm{v}^2\mid\mathcal F_k]
   \le 4\sigma_4^2\alpha_k^2\norm{X_k}^2,\\
  \E[\norm{v}^4\mid\mathcal F_k]
  &\le \sigma_4^4\alpha_k^4.
\end{align*}
For the mixed cubic term, first use Cauchy's inequality and Jensen's inequality:
\[
  \left|\E[4\ip{u}{v}\norm{v}^2\mid\mathcal F_k]\right|
  \le 4\norm{u}\E[\norm{v}^3\mid\mathcal F_k]
  \le 4\sigma_4^3\alpha_k^3\norm{X_k}.
\]
We now apply Young's inequality in the form
\[
  c_1\alpha_k^3 t
  \le \frac{\lambda}{4}\alpha_k t^4+C\alpha_k^{11/3},
  \qquad t\ge0, \label{eq:young-fourth}
\]
which follows from $ab\le \varepsilon a^4 + C\varepsilon^{-1/3}b^{4/3}$ by taking $a=t$, $b=c_1\alpha_k^3$, and $\varepsilon=(\lambda/4)\alpha_k$. Since $\alpha_k\le1$, $\alpha_k^{11/3}\le\alpha_k^3$. Therefore
\begin{equation}
  \left|\E[4\ip{u}{v}\norm{v}^2\mid\mathcal F_k]\right|
  \le \frac{\lambda}{4}\alpha_k\norm{X_k}^4+C\alpha_k^3. \label{eq:mixed-cubic}
\end{equation}
Finally,
\[
  \norm{u}^4\le(1-\lambda\alpha_k)^4\norm{X_k}^4.
\]
For $\lambda\alpha_k\le1/4$, the elementary inequality
\[
  (1-x)^4\le 1-\frac{15}{8}x,\qquad 0\le x\le\frac14,
\]
with $x=\lambda\alpha_k$ gives
\[
  \norm{u}^4\le \left(1-\frac{15\lambda}{8}\alpha_k\right)\norm{X_k}^4.
\]
Combining this estimate with \eqref{eq:mixed-cubic} and the remaining second- and fourth-order noise terms gives
\[
  \E[\norm{X_{k+1}}^4\mid\mathcal F_k]
  \le \left(1-\frac{13\lambda}{8}\alpha_k\right)\norm{X_k}^4
       +C\alpha_k^2\norm{X_k}^2+C\alpha_k^3.
\]
Absorbing constants and taking expectation,
\begin{equation}
  n_{k+1}\le(1-c_4\alpha_k)n_k+C\alpha_k^2m_k+C\alpha_k^3,
  \qquad n_k\coloneqq\E\norm{X_k}^4, \label{eq:n-recursion-pre}
\end{equation}
for $c_4=13\lambda/8$ replaced by any smaller positive constant. The already-proved second moment bound gives $m_k\le C\alpha_k$, hence
\begin{equation}
  n_{k+1}\le(1-c_4\alpha_k)n_k+C\alpha_k^3. \label{eq:n-recursion}
\end{equation}
We prove $n_k\le C_4\alpha_k^2$ by induction. Choose $C_4$ so large that $n_0\le C_4(\alpha_0 k_0^{-a})^2$. Since
\[
  \frac{\alpha_{k+1}^2}{\alpha_k^2}=\left(1+\frac1{\bar k}\right)^{-2a}
  \ge1-\frac{2a}{\bar k},
\]
and since $\bar k^{-1}=o(\alpha_k)$, after increasing $k_0$ if necessary,
\begin{equation}
  \alpha_{k+1}^2\ge \alpha_k^2\left(1-\frac{c_4}{2}\alpha_k\right). \label{eq:alpha2-compare}
\end{equation}
Assume $n_k\le C_4\alpha_k^2$. From \eqref{eq:n-recursion},
\[
  n_{k+1}\le C_4\alpha_k^2(1-c_4\alpha_k)+C\alpha_k^3
  =C_4\alpha_k^2\left(1-c_4\alpha_k+\frac{C}{C_4}\alpha_k\right).
\]
Choosing $C_4\ge 2C/c_4$ gives
\[
  n_{k+1}\le C_4\alpha_k^2\left(1-\frac{c_4}{2}\alpha_k\right)
  \le C_4\alpha_{k+1}^2.
\]
This completes the induction and proves the fourth moment bound.

\paragraph{Interpolation.}
Let $p=2+2\rho\in[2,4]$. By Lyapunov's inequality,
\[
  \E\norm{X_k}^p\le \left(\E\norm{X_k}^4\right)^{p/4}
  \le C\left(\alpha_k^2\right)^{p/4}
  =C\alpha_k^{p/2}
  =C\alpha_k^{1+\rho}.
\]
This proves \eqref{eq:fast-interp} and the lemma.
\end{proof}

\end{document}
